% !TeX root = ../main.tex
\begin{frame}{课程评价——臧笛名}
  \begin{SlideCanvas}
    \node[anchor=north west,inner sep=0pt,text width=850pt,align=left,
      text=BICNavy] at (55,80) {\begin{minipage}{850pt}
      \usebeamerfont{normal text}\raggedright
      \setlength{\parindent}{0pt}\setlength{\parskip}{4pt}
      \setlength{\columnsep}{26pt}\setlength{\multicolsep}{0pt}
      \clubpenalty=10000\widowpenalty=10000
{\usebeamerfont{BIC point}课程学习体会\par}\nobreak
通过这门课程的学习，我了解到了医疗传感器、核磁共振成像、人工智能及MEMS器件等不同方向的知识，开阔了眼界，也学到了不少平时接触较少的内容。课程内容比较丰富，但部分知识专业性较强，理解起来有一定难度。\par
\par\addvspace{22pt}{\usebeamerfont{BIC point}意见与建议\par}\nobreak
建议今后适当增加基础知识的介绍，多结合实际案例进行讲解，并适当放慢重点内容的讲解速度。同时，可以增加课堂交流的机会，多介绍相关技术在实际生活中的应用，让课程更加生动有趣，也能帮助我们更好地理解所学知识。\par
    \end{minipage}};
  \end{SlideCanvas}
  \note{原始材料为用户提供的四份课程反馈之一，正文逐字保留，仅调整标题、段落和编号排版。臧笛名《医工交叉研究前沿进展》课程反馈与建议。通过这门课程的学习，我了解到了医疗传感器、核磁共振成像、人工智能及MEMS器件等不同方向的知识，开阔了眼界，也学到了不少平时接触较少的内容。课程内容比较丰富，但部分知识专业性较强，理解起来有一定难度。\par 建议今后适当增加基础知识的介绍，多结合实际案例进行讲解，并适当放慢重点内容的讲解速度。同时，可以增加课堂交流的机会，多介绍相关技术在实际生活中的应用，让课程更加生动有趣，也能帮助我们更好地理解所学知识。}
\end{frame}
